% ID: SEL-D0080-OB-E01
% ESTADO: cotejada. Segunda pasada visual y borrador compilado revisados.
% Fuente DOC: Selectividad/sel_2012_matematicas_aplicadas/2 MATEMÁT. APLIC. CC. SS/EXAMEN 5-ANDALUCÍA.doc
% SHA-256 original: 87b729cdd4632386690a9390bdfa17dd7d5c902a2cf7e5801feea45e2b181969
% Página convertida: 2; opción B; ejercicio 1.

Los alumnos de 2º de Bachillerato organizan una venta de pasteles para el viaje de fin de curso. Venden pasteles grandes, que necesitan 2 huevos, 5 terrones de azúcar y 100 g de harina cada uno, y pasteles pequeños, que necesitan 1 huevo, 3 terrones de azúcar y 80 g de harina cada uno.
\begin{enumerate}[label=\alph*)]
\item \textbf{(0.5 puntos)} Presente en una matriz $M$, de dimensión $3\mathrm{x}2$, las cantidades de los elementos necesarios para la elaboración de un pastel grande y uno pequeño.
\item \textbf{(0.5 puntos)} Si desean fabricar 20 pasteles de una clase y 30 de otra, escriba las dos matrices columna, $A$ (20 grandes y 30 pequeños) y $B$ (30 grandes y 20 pequeños) que representan este reparto.
\item \textbf{(1.5 puntos)} Calcule los productos $M\cdot A$ y $M\cdot B$ e indique si con 8 docenas de huevos, 200 terrones de azúcar y 5 kg de harina se pueden elaborar 20 pasteles grandes y 30 pequeños. \textquestiondown{}Y 30 grandes y 20 pequeños?
\end{enumerate}
